\section{Event inference}

Algorithm~\ref{alg:infer} (Algorithm~1 of the paper) replaces thresholding and
median filtering. Inputs are the presence
probabilities $\hat p^{\pres}_{c,t}$, proposals $\hat r_{c,t}$, a maximum
$k$ events per class (default 15), and the $m$ most active classes (default all).
For each selected class:
\begin{enumerate}
\item sort frames by presence;
\item take the unprocessed frame $t$ with the highest presence; its proposal
      $\hat r_{c,t}$ is the candidate;
\item set its event score $\sigma$ to the mean presence over the proposal;
\item emit $(c,\;t-\hat d^{\on},\;t+\hat d^{\off},\;\sigma)$;
\item suppress every proposal overlapping the chosen interval;
\item repeat until $k$ events or no proposals remain.
\end{enumerate}
This is a non-maximum-suppression-like procedure. It is not differentiable; it
only converts learned proposals into a final event set.

The baseline is the usual frame-wise post-processing:
\begin{equation}
\tilde p_t
=
\median\left(p_{t-\lfloor L/2\rfloor},\ldots,p_{t+\lfloor L/2\rfloor}\right),
\end{equation}
followed by threshold $\tau$ and grouping of consecutive active frames.

\begin{algorithm}[t]
\caption{Event inference (Algorithm~1 of the paper)}
\label{alg:infer}
\begin{algorithmic}[1]
\Require presence $\hat p^{\pres}_{c,t}$; proposals $\hat r_{c,t}$;
         max events per class $k$; max classes $m$
\Ensure event set $\mathcal{E}$
\State $s_c \gets \operatorname{mean}_t \hat p^{\pres}_{c,t}$ for every class
\State select the $m$ classes with the largest $s_c$
\State $\mathcal{E} \gets \emptyset$
\For{each selected class $c$}
  \State order frames by $\hat p^{\pres}_{c,t}$ descending
  \State $removed_t \gets \mathrm{false}$ for all $t$
  \State $\mathcal{E}_c \gets \emptyset$ \Comment{events selected for class $c$}
  \For{$t$ in that order}
    \If{$removed_t$ or $|\mathcal{E}_c| \ge k$}
      \State continue
    \EndIf
    \State $r \gets [t-\hat d^{\on}_{c,t},\,t+\hat d^{\off}_{c,t}]$,
           clipped to $[0,T_{\mathrm{ds}}]$
    \State $\sigma \gets \operatorname{mean}_{u\in r}\hat p^{\pres}_{c,u}$
    \State $\mathcal{E}_c \gets \mathcal{E}_c\cup\{(c,r,\sigma)\}$
    \State mark every proposal overlapping $r$ as removed
  \EndFor
  \State $\mathcal{E} \gets \mathcal{E}\cup\mathcal{E}_c$
\EndFor
\State \Return $\mathcal{E}$
\end{algorithmic}
\end{algorithm}

Algorithm~\ref{alg:infer} is greedy and has no learned parameters. The score
$\sigma$ is the mean presence over the chosen interval, so a proposal that
covers a confident event ranks above one that covers a weak region. Suppressing
every overlapping proposal keeps the output set small, but it also means two
same-class events that overlap in time collapse into one. The method cannot
represent same-class polyphony, and the evaluation keeps the raw annotations so
those events still count as false negatives.

The inference step is also the place where $k$ and $m$ enter. A small $k$
limits repeated detections, while a small $m$ skips classes with low mean
presence. In our runs, EPN inference and a tuned median filter were close on
the test split (43.58 vs 43.00 PSDS1 for the wide model), which suggests that
presence quality limits this reproduction more than the choice of event
extraction.

